% =============================================================================
% CONCLUSION AND PERSPECTIVES
% =============================================================================
% COMPLETELY REWRITTEN — Old Knowledge Graph / Data Fabric content removed
% Project: A Context-Aware Framework for Streaming Data Quality Monitoring
% Date: April 24, 2026
% =============================================================================

\section{Summary of Contributions}
\label{sec:conclusion-contributions}

We presented StreamDQ: a context-aware framework for streaming data quality
monitoring, implemented in Apache Flink and evaluated on NYC transportation
data streams. We summarize the four primary contributions of this work.

\textbf{C1: First streaming GPS trajectory validation on real GTFS-realtime
feeds.} We implemented three cross-record rules — CRS001 (Haversine-based GPS
speed bounds $[2.0,\ 100.0]$ km/h), CRS002 (GPS jump detection $>400$ m
in $\leq 30$ s), and CRS003 (event deduplication with 300-second window) —
as Java \texttt{KeyedProcessFunction} operators invoked via gRPC from the
Flink Python pipeline. These rules are evaluated exclusively on the NYC MTA
Bus GTFS-realtime stream, which provides GPS coordinates; they cannot be
evaluated on the NYC TLC Yellow Taxi dataset (zone-level only).
\texttt{[TIER-2\ ESTIMATED]}

\textbf{C2: Hierarchical context-aware threshold fallback (L0--L5).} We
designed a five-level hierarchical threshold system that traverses from the
finest-grained context cell (L0: hour $\times$ zone $\times$ weekday/weekend)
to physics priors (L5), with graceful degradation through L1--L4. The
minimum sample thresholds per level (L0 $\geq 100$, L1 $\geq 50$, L2 $\geq 25$,
L3 $\geq 10$, L4 $\geq 5$) are derived from power analysis (Cohen's $d = 0.30$,
$\alpha = 0.01$, power $= 0.80$). We validated that this hierarchical
fallback mechanism is absent from all surveyed frameworks (Great Expectations,
Soda Core, Stream DaQ, METER, Ada-Context, Deequ, GTFS Validator, and
7 others). Expected L0 coverage is 2--5\% of events due to zone-hour sparsity
in the NYC TLC dataset.
\texttt{[TIER-2\ ESTIMATED]}

\textbf{C3: Ground-truth evaluation methodology.} We implemented a complete
evaluation pipeline comprising synthetic anomaly injection, per-entity ground-truth
tracking in PostgreSQL, and precision/recall/F1 metrics with 95\% bootstrap
confidence intervals (1,000 iterations). We explicitly acknowledge what cannot
be measured: CRS003 recall is blocked by NG-4 (replay suppression gate), and
distributed Flink latency is not benchmarked. This honest methodology is
modeled on Exathlon \citep{exathlon2021}, which established precedent for
benchmarking streaming anomaly detection.
\texttt{[TIER-1\ VERIFIED\ for\ methodology;\ TIER-2\ ESTIMATED\ for\ results]}

\textbf{C4: ML-augmented threshold calibration.} We designed a three-component
ML pipeline — Bayesian Optimization (Priority 1, GO), Isolation Forest (Priority 2,
conditional), and XGBoost (Priority 3, conditional) — that calibrates context
thresholds without overriding rule-based violation detection. Rules remain
authoritative at all times. LSTM (Priority 4, NO-GO) is excluded due to
unconfirmed GPS training data and high redundancy with CRS002.
\texttt{[TIER-2\ ESTIMATED]}

\section{Key Findings}
\label{sec:conclusion-findings}

We report the following key findings based on the design and methodology
established in this thesis.

\textbf{CRS001 and CRS002 are designed as specified with physics-based
thresholds.} CRS001 uses Haversine distance with the $[2.0,\ 100.0]$ km/h
bound calibrated for NYC MTA Bus urban operations. The lower bound of
2.0 km/h excludes GPS jitter; the upper bound of 100.0 km/h closes the
detection gap for moderate spoofing (20--100 km/h). CRS002 uses the
400-meter/30-second threshold calibrated for the 30-second GTFS-realtime
update interval. Both rules use L5 physics priors exclusively; they have
no context dependency.
\texttt{[TIER-2\ ESTIMATED]}

\textbf{Context-aware thresholds provide value at L0 cells but coverage is
sparse.} The expected L0 coverage of 2--5\% of events means that most
events are evaluated at coarser fallback levels (L3 or L4). At L4 (global
fallback), the threshold is identical to a static threshold, so no
context-specific improvement is available. The aggregate $\Delta$F1 across
all events is expected at 1--2 pp, not the 5 pp target which applies only
to L0 cells. This finding is critical for interpreting RQ1 honestly.
\texttt{[TIER-2\ ESTIMATED]}

\textbf{L4 global fallback is necessary but dilutes context specificity.}
The hierarchical fallback ensures that every event receives a threshold
(even if only the physics prior), preventing silent failures. However, the
L4 global threshold pools all events regardless of context, which may
increase the false-negative rate for events with unusual but legitimate
values. This trade-off is inherent to the fallback design.

\textbf{ML augmentation is a measured contribution.} Phase 3 implementation
is mandatory. Whether ML provides measurable F1 improvement over rule-only
thresholds is an empirical question that requires benchmark execution.
Positive, negative, and mixed results are all scientifically valid and
will be reported as such.

\textbf{TQS correlation with downstream quality remains to be validated.}
The non-circular TQS design (five stream-intrinsic dimensions, weighted
sum) distinguishes StreamDQ from prior work that defines TQS as
$1 - \text{violation\_rate}$. Whether TQS meaningfully predicts downstream
quality (e.g., ETA prediction accuracy) requires independent measurement
using ETA error computed from GTFS static schedules and published delays.
\texttt{[TIER-2\ ESTIMATED]}

\section{Future Work}
\label{sec:conclusion-future}

Five directions for future work emerge from this research.

\textbf{1. Fix NG-4: replay suppression gate for CRS003 evaluation.}
The current code tags synthetic duplicate events as \texttt{is\_replay=True},
which triggers the replay suppression gate and prevents violation emission.
Fixing this bug requires tagging synthetic duplicates as \texttt{is\_replay=False}.
Once fixed, CRS003 recall can be measured alongside precision, completing
the CRS rule evaluation.

\textbf{2. Implement D4: External context (holiday indicator).} The
five-dimensional context decomposition currently operates on four dimensions
(D1, D2, D3, D5). Implementing the holiday indicator in D4 would upgrade
the context model to full five-dimensional operation. Holidays affect
both taxi fare distributions (surcharges, airport trips) and bus ridership
patterns. A holiday lookup table for NYC and integration with the context
extraction pipeline are the primary implementation steps.

\textbf{3. Deploy Flink cluster for distributed benchmarks.} LocalPipeline
results are not equivalent to distributed Flink performance. Deploying a
multi-node Flink cluster with Kafka partitioning and RocksDB state would
enable proper latency and throughput measurements. This is necessary for
claiming P99 latency numbers that are representative of real-world
streaming deployment.

\textbf{4. Re-evaluate LSTM if GPS training data becomes available.}
The LSTM trajectory prediction component is explicitly excluded (Priority 4,
NO-GO) due to unconfirmed 6-month historical NYC MTA Bus GPS archive.
If this archive is confirmed to exist, LSTM should be re-evaluated under
two conditions: (a) CRS002 recall falls below 60\% on real GPS anomalies,
and (b) LSTM provides measurable improvement over CRS002 alone. LSTM
would serve only as an alert elevation mechanism, never as authoritative
validation.

\textbf{5. Extend CRS rules to other GTFS feeds beyond NYC.} The
CRS001 and CRS002 rules are calibrated for NYC MTA Bus operations (urban
bus, 30-second update interval, 100 km/h speed limit). Generalizing to
other transit agencies requires parameter recalibration for different
vehicle types (subway, rail, ferry), geographic contexts (rural vs. urban),
and update frequencies. A calibration framework that automates threshold
derivation from the GTFS-realtime feed characteristics would enable
broader deployment.

\section*{Chapter Summary}
\label{sec:conclusion-summary}

This thesis presented StreamDQ: a context-aware framework for streaming
data quality monitoring, designed for GPS trajectory validation on real
GTFS-realtime transit feeds. We implemented three layers of data quality
rules (syntactic, semantic, cross-record), a hierarchical context-aware
threshold system with L0--L5 fallback, a non-circular Trajectory Quality
Scoring mechanism, and ML-augmented threshold calibration. We designed a
rigorous evaluation methodology with synthetic anomaly injection, ground-truth
tracking, and bootstrap confidence intervals. We explicitly acknowledged
seven limitations, including the NG-4 blocker for CRS003 recall, the sparse
L0 coverage, and the distinction between LocalPipeline and distributed Flink
benchmarks. StreamDQ is a research and education platform that demonstrates
the feasibility and measurability of streaming GPS trajectory validation
for public transit data quality monitoring.
